% Generated draft: requires mathematical and source review.
\documentclass[11pt]{article}
\usepackage[margin=25mm]{geometry}
\usepackage{fontspec}
\setmainfont{Latin Modern Roman}
\usepackage{amsmath,amssymb,mathtools,unicode-math}
\setmathfont{Latin Modern Math}
\usepackage{enumitem,needspace,xurl,hyperref}
\setlength{\parindent}{0pt}
\setlength{\parskip}{5pt}
\setlength{\emergencystretch}{3em}
\newcommand{\N}{\mathbb{N}}
\newcommand{\Z}{\mathbb{Z}}
\newcommand{\Q}{\mathbb{Q}}
\newcommand{\R}{\mathbb{R}}
\newcommand{\C}{\mathbb{C}}
\newcommand{\F}{\mathbb{F}}
\newcommand{\A}{\mathbb{A}}
\newcommand{\PP}{\mathbb{P}}
\newcommand{\E}{\mathbb{E}}
\newcommand{\eps}{\varepsilon}
\newcommand{\dd}{\,\mathrm{d}}
\begin{document}
\section*{UnsolvedMath ID: 30003633}
\textbf{Source problem code:} catalogue entry
\subsection{Definitions and mathematical statement}
Let \(m\in\Z\) with \(m\geq2\), and let \(n\in\Z_{\geq0}\). Fix distinct marked points \(p_1,\ldots,p_m\) in the interior of the disk \(D^2\). An \(m\)-component string link consists of \(m\) disjoint, ordered, oriented strands in \(D^2\times[0,1]\), with strand \(i\) running from \((p_i,0)\) to \((p_i,1)\), and with the usual fixed behavior near the boundary. String links modulo concordance relative to the boundary form a group \(\mathcal C_m\) under stacking. The supplied record does not specify whether concordance is smooth or locally flat topological; the intended category must be checked in the cited source.

Write \(\mathcal F_n^m\) and \(\mathcal F_{n+\frac12}^m\) for the integer and subsequent half-integer terms of the \(n\)-solvable filtration of \(\mathcal C_m\) referred to in the source. Thus the notation \(\mathcal F_{n.5}^m\) in the record means \(\mathcal F_{n+\frac12}^m\). These symbols retain the source's intended solvability convention: the supplied record does not give the defining four-dimensional solvability conditions, so a complete intrinsic definition of these subgroups cannot be recovered from the supplied text alone. The quotient notation presupposes \(\mathcal F_{n+\frac12}^m\trianglelefteq\mathcal F_n^m\).

\textbf{Question.} With these source-dependent conventions fixed, is it true that
\[
\forall m\in\Z_{\geq2}\;\forall n\in\Z_{\geq0},\qquad
\mathcal F_n^m/\mathcal F_{n+\frac12}^m
\text{ is nonabelian}?
\]
Equivalently, writing \([x,y]=xyx^{-1}y^{-1}\), the question asks whether
\[
\forall m\in\Z_{\geq2}\;\forall n\in\Z_{\geq0}\;
\exists x,y\in\mathcal F_n^m
\quad [x,y]\notin\mathcal F_{n+\frac12}^m.
\]
The intended concordance category and precise solvability convention remain unresolved qualifications of this statement; the displayed question does not assert that the answer is affirmative.
\subsection{Short English statement}
For every integer number of components \(m\geq2\) and every integer level \(n\geq0\), must there be two \(n\)-solvable string-link concordance classes whose commutator does not lie in the next half-level of the solvable filtration? The record leaves the concordance category and technical definition of solvability unspecified. Its supplied status is partially solved, subject to the dated source qualification below.
\subsection{Sources}
\begin{itemize}
\item Local UnsolvedMath registry, permanent ID \(30003633\), title \emph{Nonabelian Quotients of the String Link Solvable Filtration}, exact source problem code \texttt{catalogue entry}; collection: \emph{Oberwolfach Reports Open Problems}. No individual proposer is identified in the supplied record; its proposed year is 2017.
\item Supplied source citation: \emph{Mini-Workshop: Interactions between Low-dimensional Topology and Complex Algebraic Geometry (2018)}, Oberwolfach Reports, pp.\ 2987--3033, DOI \texttt{10.4171/owr/2017/50}. \url{https://doi.org/10.4171/owr/2017/50}. The supplied title gives 2018, whereas the report identifier and proposed-year field indicate 2017; this discrepancy is unresolved.
\item Supplied literature reference labeled ``Otto'': \url{https://msp.org/agt/2014/14-5/agt-v14-n5-p05-s.pdf}. No further bibliographic details or precise theorem references are supplied.
\item Supplied literature reference labeled ``Harvey'': \url{https://arxiv.org/abs/math/0609378}. No further bibliographic details or precise theorem references are supplied.
\item Supplied literature reference labeled ``Oberwolfach Report 50/2017'': \url{https://doi.org/10.4171/owr/2017/50}.
\item Literature assessment dated exactly \texttt{2026-08-21}: the supplied curation labels the problem partially solved and reports substantial nontrivial and nonabelian successive quotients at low levels and for suitable numbers of components, while describing the full extracted nonabelianity question as unresolved in general. This is supplied dated curation derived from the OWR record, not an independently verified literature assessment. The record explicitly calls for independent verification of time-sensitive status and cited claims.
\item The registry labels the statement \texttt{corrected\_verified} and explains that mathematical content was preserved while notation, typography, and sentence structure were repaired. This is a supplied editorial verification label, not independent mathematical review of this draft.
\end{itemize}
\end{document}
